% Baccalauréat malgache, série D, 2011 : énoncé.
\TCDBacAttribution{D}{2011}
\providecommand{\Courbe}{(\mathscr{C})}
\providecommand{\D}{(\mathscr{D})}
\providecommand{\Plan}{(\mathscr{P})}
\providecommand{\Ouv}{(O ; \vec{u} , \vec{v})}
%\providecommand{\Oij}{(O ; \vec{\imath} , \vec{\jmath})}
\providecommand{\cm}{~\mathrm{cm}}
\providecommand{\dx}{~\mathrm{d}x}

\noindent \underline{\textsc{Exercice I}} (5 points) \newline
      Soit l'\'equation 
      \begin{equation}
            \TCDtag{E}
                  z^3 + (1-i) z^2 + (-8+4i)z - 4 - 28 i = 0 \;.
            \label{eq: Exo 1 Bac D 2011}
      \end{equation}
      \begin{enumerate}
            \item Montrer que \eqref{eq: Exo 1 Bac D 2011} admet une solution imaginaire pure $\alpha \, i$ 
                  o\`u $\alpha$ est un nombre r\'eel que l'on d\'eterminera.
            \item \begin{enumerate}
                        \item Montrer que \eqref{eq: Exo 1 Bac D 2011} peut s'\'ecrire sous la forme :
                              $(z-\alpha \, i) ( a z^2 + b z + c) = 0$ o\`u $a$, $b$ et $c$ sont des nombres 
                              complexes \`a d\'eterminer.
                        \item R\'esoudre dans $\mathbb{C}$ l'\'equation \eqref{eq: Exo 1 Bac D 2011}.
                  \end{enumerate}
            \item Dans le plan $\Plan$ muni d'un rep\`ere orthonorm\'e direct $\Ouv$, d'unit\'e $1\cm$, 
                  on donne les points $A$, $B$ et $C$ d'affixes respectives $-2i$, $3+i$ et $-4+2i$.
                  \begin{enumerate}
                        \item Placer les points $A$, $B$ et $C$ dans le plan $\Plan$.
                        \item On pose $Z = \frac{z_B - z_A}{z_C - z_A}$.
                              Ecrire $Z$ sous forme trigonom\'etrique.
                        \item En d\'eduire la nature du triangle $ABC$.
                        \item Trouver l'affixe du point $D$ pour que le quadrilat\`ere $ABCD$ soit un parall\'elogramme.
                  \end{enumerate}
            \item Soit $S$ la transformation d'expression complexe $z' = \frac{4}{3}i z - \frac{8}{3} - 2i$.
                  \begin{enumerate}
                        \item Quelle est la nature de $S$ et donner ses \'el\'ements caract\'eristiques.
                        \item Construire l'image $A'B'C'D'$ de $ABCD$ par $S^{-1}$.
                  \end{enumerate}
      \end{enumerate}
%%%%%%%%%%%%%%%%%%%%%%%%%%%%%%%%%%%%%%%%%%%%%%%%%%%%%%%%%%%%%%%%%%%%%%%%%%%%%%%%%%%%%%%%%%%%%%%
%%%%%%%%%%%%%%%%%%%                   %%%%%%%%%%%%%%%%%%%                   %%%%%%%%%%%%%%%%%%%
%%%%%%%%%%%%%%%%%%%%%%%%%%%%%%%%%%%%%%%%%%%%%%%%%%%%%%%%%%%%%%%%%%%%%%%%%%%%%%%%%%%%%%%%%%%%%%%
\noindent \underline{\textsc{Exercice II}} (5 points) 
      \begin{description}
            \item[I.] Une urne contient 8 boules indiscernables au toucher dont :
                  \begin{itemize}
                        \item 3 rouges num\'erot\'ees : 2, 2, 5,
                        \item 5 blanches num\'erot\'ees : 1, 2, 3, 4, 5.
                  \end{itemize}
                  \begin{enumerate}
                        \item Le jeu consiste \`a tirer au hasard et simultan\'ement deux boules de l'urne.
                              Calculer la probabilit\'e de chacun des \'ev\`enements suivants :
                              \begin{itemize}
                                    \item $A$ : << Obtenir deux boules de couleurs diff\'erentes >>.
                                    \item $B$ : << Obtenir deux boules de num\'eros pairs >>.
                                    \item $C$ : << Obtenir deux boules dont le produit des num\'eros est \'egal \`a 4 >>.
                              \end{itemize}
                        \item On tire au hasard et successivement, avec remise, 3 boules de l'urne.
                              Soit $X$ la variable al\'eatoire d\'efinie par le nombre de boules portant le num\'ero 5 obtenues.
                              \begin{enumerate}
                                    \item D\'eterminer l'univers image de $X$.
                                    \item Donner la loi de probabilit\'e de $X$.
                              \end{enumerate}
                  \end{enumerate}
                  \textbf{NB : On donnera les r\'esultats sous forme de fraction irr\'eductible.}
            \item[II.] On donne sur le tableau ci-dessous le nombre d'\'el\`eves d'un lyc\'ee ayant r\'eussi le 
                  Baccalaur\'eat durant 4 ann\'ees successives : 
                  \par\smallskip\noindent
                  \adjustbox{max width=\linewidth,center}{%
                  \begin{tabular}{|l|c|c|c|c|}
                  \hline
                  Ann\'ee                                 & 2007 & 2008 & 2009 & 2010 \\ \hline
                  Rang de l'ann\'ee $(x_i)$               & 1    & 2    & 3    & 4    \\ \hline
                  Nombre d'\'el\`eves en centaines $(y_i)$ & 3    & 5    & 6    & 9    \\ \hline
                  \end{tabular}%
                  }
                  \par\smallskip
                  \begin{enumerate}
                        \item Repr\'esenter le nuage des points $M_i (x_i ; y_i)_{1 \leq i \leq 4}$ associ\'e \`a cette
                              s\'erie statistique. \newline
                              Echelle : 
                                    \begin{itemize}
                                          \item sur l'axe des abscisses, prendre $1\cm$ pour repr\'esenter une unit\'e,
                                          \item sur l'axe des ordonn\'ees, placer 2 \`a l'origine des axes puis prendre
                                                $1\cm$ pour repr\'esenter $100$ \'el\`eves.
                                    \end{itemize}
                        \item D\'eterminer le point moyen $G$.
                        \item Calculer le coefficient de corr\'elation $r$ et interpr\'eter.
                        \item Ecrire l'\'equation de la droite de r\'egression de $y$ en $x$.
                        \item Combien de r\'eussites peut-on esp\'erer en 2014\,?
                  \end{enumerate}
                  \textbf{NB : Les r\'esultats seront donn\'es \`a $10^{-2}$ pr\`es.}
      \end{description}
%%%%%%%%%%%%%%%%%%%%%%%%%%%%%%%%%%%%%%%%%%%%%%%%%%%%%%%%%%%%%%%%%%%%%%%%%%%%%%%%%%%%%%%%%%%%%%%
%%%%%%%%%%%%%%%%%%%                   %%%%%%%%%%%%%%%%%%%                   %%%%%%%%%%%%%%%%%%%
%%%%%%%%%%%%%%%%%%%%%%%%%%%%%%%%%%%%%%%%%%%%%%%%%%%%%%%%%%%%%%%%%%%%%%%%%%%%%%%%%%%%%%%%%%%%%%%
\noindent \underline{\textsc{Probleme}} \newline
      On consid\`ere la fonction num\'erique $f$ d\'efinie sur $\mathbb{R}$ par : $f(x) = 2x + 1 - x \expo{-x}$.
      On note $\Courbe$ sa courbe repr\'esentative dans un plan muni d'un rep\`ere orthonorm\'e \Oij{} d'unit\'e $2\cm$.
      \begin{enumerate}
            \item Soit $g$ la fonction num\'erique d\'efinie sur $\mathbb{R}$ par : $g(x) = (x-1) \expo{-x} + 2$.
                  \begin{enumerate}
                        \item Etudier les variations de $g$ et dresser son tableau de variation.
                        \item D\'emontrer que l'\'equation $g(x) = 0$ admet une solution unique not\'ee $\alpha$ 
                              dans $\mathbb{R}$ et que $-1 < \alpha < 0$.
                        \item En d\'eduire le signe de $g(x)$ suivant les valeurs de $x$.
                  \end{enumerate}
            \item \begin{enumerate}
                        \item Calculer $\lim_{x\rightarrow -\infty} f(x) $ et $\lim_{x\rightarrow +\infty} f(x)$.
                        \item Calculer la d\'eriv\'ee $f'(x)$ et montrer que : $f'(x) = g(x)$ pour tout $x$ \'el\'ement de 
                              $\mathbb{R}$.
                        \item Montrer que $f(\alpha ) = 1 + \frac{2 \alpha^2}{\alpha - 1}$.
                        \item Dresser le tableau de variation de $f$.
                  \end{enumerate}
            \item \begin{enumerate}
                        \item Calculer $\lim_{x\rightarrow -\infty} \frac{f(x)}{x}$ et interpr\'eter graphiquement ce r\'esultat.
                        \item D\'emontrer que la droite $\D : y = 2x + 1$ est asymptote \`a $\Courbe$ au voisinage de $+\infty$.
                        \item Etudier la position de $\Courbe$ par rapport \`a $\D$.
                  \end{enumerate}
            \item \begin{enumerate}
                        \item Ecrire une \'equation de la tangente $(T)$ \`a $\Courbe$ au point d'abscisse $x_0 = 0$.
                        \item Montrer qu'il existe un point $A$ de $\Courbe$ o\`u $\Courbe$ admet une 
                              tangente $(T')$ parall\`ele \`a $\D$. Trouver les coordonn\'ees du point $A$.
                  \end{enumerate}
            \item Tracer $\D$, $(T)$, $(T')$ et $\Courbe$ dans le m\^eme rep\`ere. 
                  Pour la construction, on prendra $\alpha = 0,4$ et $f(\alpha) = 0,8$.
            \item \begin{enumerate}
                        \item A l'aide d'une int\'egration par partie ; 
                              calculer $I_{\lambda} = \int_{0}^{\lambda} x \expo{-x} \dx$.
                        \item D\'eterminer, en $\cm ^2$, l'aire $\mathscr{A}(\lambda)$ du domaine plan limit\'e par 
                              $\Courbe$, $\D$ et les droites d'\'equations $x = 0$ et $x = \lambda$ ($\lambda > 0$).
                        \item Calculer $\lim_{\lambda \rightarrow +\infty} \mathscr{A}(\lambda)$.
                  \end{enumerate}
      \end{enumerate}





